\documentclass[12pt, a4paper]{article}

\usepackage{amsmath, amssymb, amsthm}
\usepackage{geometry}
\usepackage{hyperref}
\usepackage{listings}
\usepackage{xcolor}
\usepackage{titlesec}
\usepackage{parskip}
\usepackage{booktabs}
\usepackage{graphicx}
\usepackage{fancyhdr}

\geometry{margin=2.5cm}

\hypersetup{
    colorlinks=true,
    linkcolor=blue!70!black,
    urlcolor=blue!70!black
}

\lstset{
    language=,
    basicstyle=\ttfamily\small,
    backgroundcolor=\color{gray!10},
    frame=single,
    rulecolor=\color{gray!40},
    breaklines=true,
    keywordstyle=\color{blue!70!black}\bfseries,
    commentstyle=\color{green!50!black},
    stringstyle=\color{red!70!black},
}

\theoremstyle{definition}
\newtheorem{theorem}{Theorem}[section]
\newtheorem{definition}[theorem]{Definition}
\newtheorem{corollary}[theorem]{Corollary}
\newtheorem{remark}[theorem]{Remark}

\pagestyle{fancy}
\fancyhf{}
\rhead{Dean Foulds}
\lhead{Aircraft Parts Volatility — Formal Proofs}
\cfoot{\thepage}

\title{
    \vspace{-1cm}
    {\Large\bfseries Formal Mathematical Proofs for the}\\[0.4em]
    {\LARGE\bfseries Aircraft Parts Volatility Index (APVI)}\\[0.6em]
    {\large A Lean 4 \& Mathlib Verification}
}
\author{
    Dean Foulds\\
    \small Data Scientist \& ML Engineer\\
    \small \href{https://deanfoulds.xyz}{deanfoulds.xyz} \quad
    \href{https://github.com/Dean-Foulds/lean_projects}{github.com/Dean-Foulds/lean\_projects}
}
\date{\today}

\begin{document}

\maketitle
\tableofcontents
\newpage

% ---------------------------------------------------------------
\section{Introduction}
% ---------------------------------------------------------------

This document presents formal mathematical proofs underpinning the
\textbf{Aircraft Parts Volatility Index (APVI)}, a composite risk metric
that combines equity market volatility (via the XAR ETF) with Producer Price
Index (PPI) volatility for the aerospace parts sector.

The proofs are machine-verified using \textbf{Lean 4} and the
\textbf{Mathlib4} library — one of the largest formal mathematics libraries
in existence. Lean verification guarantees that each theorem holds under
the stated hypotheses with no logical gaps.

The APVI is defined as:
\[
    \text{APVI} = \frac{\sigma_{\text{equity}} + \sigma_{\text{PPI}}}{2}
\]

where $\sigma_{\text{equity}}$ is the 90-day realised volatility of the XAR
aerospace ETF and $\sigma_{\text{PPI}}$ is the rolling 12-month standard
deviation of log PPI returns for aircraft parts (FRED series
\texttt{PCU336413336413P}).

% ---------------------------------------------------------------
\section{Definitions}
% ---------------------------------------------------------------

\begin{definition}[Volatility Spike]
Let $v, \mu, \sigma, k \in \mathbb{R}$ with $\sigma > 0$. A value $v$ is
said to constitute a \emph{volatility spike} at threshold $k$ if:
\[
    v > \mu + k\sigma
\]
In Lean 4 this is encoded as:
\begin{lstlisting}
def is_spike (v μ σ k : ℝ) : Prop := v > μ + k * σ
\end{lstlisting}
This definition captures the standard statistical notion of an outlier:
a value exceeding the mean by more than $k$ standard deviations. In the
context of global supply chain disruption events (e.g.\ the Suez Canal
blockage of 2021 or the COVID-19 pandemic of 2020), $v$ represents
observed volatility and $\mu, \sigma$ are computed over a historical
baseline window.
\end{definition}

\begin{definition}[Realised Volatility]
For a sample of $n$ log returns $r_1, \ldots, r_n$, the annualised
realised volatility is:
\[
    \sigma_{\text{realised}} = \sqrt{\frac{1}{n}\sum_{i=1}^{n} r_i^2}
    \times \sqrt{252}
\]
In the APVI model, a 90-day rolling window is used ($n = 90$).
\end{definition}

\begin{definition}[Composite APVI]
Given equity volatility $\sigma_e \geq 0$ and PPI volatility
$\sigma_p \geq 0$, the composite index is:
\[
    \text{APVI} = \frac{\sigma_e + \sigma_p}{2}
\]
\end{definition}

% ---------------------------------------------------------------
\section{Theorems and Proofs}
% ---------------------------------------------------------------

\subsection{Spike Detection is Monotone in the Threshold}

\begin{theorem}[Monotonicity of Spike Detection]
\label{thm:spike_monotone}
Let $v, \mu, \sigma \in \mathbb{R}$ with $\sigma > 0$, and let
$k_1 < k_2$. If $v$ is a spike at threshold $k_2$, then $v$ is also
a spike at the lower threshold $k_1$:
\[
    k_1 < k_2 \;\wedge\; v > \mu + k_2\sigma
    \;\Longrightarrow\; v > \mu + k_1\sigma
\]
\end{theorem}

\begin{proof}
Since $\sigma > 0$ and $k_1 < k_2$, we have $k_1\sigma < k_2\sigma$,
hence $\mu + k_1\sigma < \mu + k_2\sigma < v$. \qed
\end{proof}

\begin{remark}
This theorem formalises the intuitive property that a stricter threshold
($k_2 > k_1$) is harder to trigger. For the APVI model, this means that
a spike flagged at the 99th percentile ($k \approx 2.33$) is automatically
flagged at the 95th percentile ($k \approx 1.65$) as well. Lean proof:
\begin{lstlisting}
theorem spike_monotone_k (v μ σ k₁ k₂ : ℝ)
    (hσ : 0 < σ) (hk : k₁ < k₂)
    (h : is_spike v μ σ k₂) : is_spike v μ σ k₁ := by
  unfold is_spike at *
  linarith [mul_lt_mul_of_pos_right hk hσ]
\end{lstlisting}
\end{remark}

\subsection{Non-Negativity of Realised Volatility}

\begin{theorem}[Realised Volatility is Non-Negative]
\label{thm:vol_nonneg}
For any sequence of returns $r : \{0,\ldots,n-1\} \to \mathbb{R}$:
\[
    \sqrt{\sum_{i=0}^{n-1} \frac{r_i^2}{n}} \;\geq\; 0
\]
\end{theorem}

\begin{proof}
The expression under the square root is a sum of non-negative terms
divided by a positive integer, hence non-negative. The square root of
a non-negative real is non-negative by definition. \qed
\end{proof}

\begin{remark}
This is a fundamental sanity property of any volatility measure.
Lean proof:
\begin{lstlisting}
theorem realized_vol_nonneg (n : ℕ) (returns : Fin n → ℝ) :
    0 ≤ Real.sqrt (∑ i, (returns i) ^ 2 / n) := Real.sqrt_nonneg _
\end{lstlisting}
\end{remark}

\subsection{Non-Negativity of the APVI}

\begin{theorem}[APVI is Non-Negative]
\label{thm:apvi_nonneg}
If $\sigma_e \geq 0$ and $\sigma_p \geq 0$, then:
\[
    \frac{\sigma_e + \sigma_p}{2} \;\geq\; 0
\]
\end{theorem}

\begin{proof}
Follows immediately from the non-negativity of addition and division
by a positive constant. \qed
\end{proof}

\begin{remark}
Lean proof:
\begin{lstlisting}
theorem apvi_nonneg (equity_vol ppi_vol : ℝ)
    (h1 : 0 ≤ equity_vol) (h2 : 0 ≤ ppi_vol) :
    0 ≤ (equity_vol + ppi_vol) / 2 := by linarith
\end{lstlisting}
\end{remark}

\subsection{APVI is Bounded by its Components}

\begin{theorem}[APVI Upper Bound]
\label{thm:apvi_bound}
For any $\sigma_e, \sigma_p \in \mathbb{R}$:
\[
    \frac{\sigma_e + \sigma_p}{2} \;\leq\; \max(\sigma_e, \sigma_p)
\]
\end{theorem}

\begin{proof}
Without loss of generality, assume $\sigma_e \geq \sigma_p$, so
$\max(\sigma_e, \sigma_p) = \sigma_e$. Then:
\[
    \frac{\sigma_e + \sigma_p}{2} \leq \frac{\sigma_e + \sigma_e}{2} = \sigma_e
\]
The case $\sigma_p > \sigma_e$ is symmetric. \qed
\end{proof}

\begin{remark}
This theorem guarantees that the composite APVI can never exceed its
most extreme component — an important property for risk interpretation.
A spike in one measure is smoothed by averaging, so both components
must be elevated for the APVI to reach extreme levels. Lean proof:
\begin{lstlisting}
theorem apvi_le_max (equity_vol ppi_vol : ℝ) :
    (equity_vol + ppi_vol) / 2 ≤ max equity_vol ppi_vol := by
  simp [max_def]
  split_ifs with h
  · linarith
  · linarith
\end{lstlisting}
\end{remark}

\subsection{Composite Spike Detection}

\begin{theorem}[Bilateral Spike Implies APVI Spike]
\label{thm:apvi_spike}
Let $\sigma_e, \sigma_p, \mu_e, \mu_p, s_e, s_p, k \in \mathbb{R}$.
If both the equity volatility and PPI volatility are spikes at
threshold $k$:
\[
    \sigma_e > \mu_e + k s_e \;\wedge\; \sigma_p > \mu_p + k s_p
\]
then the composite APVI also spikes above its composite baseline:
\[
    \frac{\sigma_e + \sigma_p}{2} > \frac{\mu_e + \mu_p}{2}
    + k\cdot\frac{s_e + s_p}{2}
\]
\end{theorem}

\begin{proof}
Adding the two spike inequalities:
\[
    \sigma_e + \sigma_p > (\mu_e + k s_e) + (\mu_p + k s_p)
    = (\mu_e + \mu_p) + k(s_e + s_p)
\]
Dividing both sides by 2 gives the result. \qed
\end{proof}

\begin{remark}
This is the key theorem linking individual market signals to the
composite index. It formally proves that simultaneous stress in both
the equity market (XAR blue line spike) and the supply chain (PPI red
line spike) — as observed during COVID-19 in Q1 2020 and the Suez Canal
blockage in Q1 2021 — necessarily produces a composite APVI spike.
Lean proof:
\begin{lstlisting}
theorem apvi_spike (eq_vol ppi_vol μ_eq μ_ppi σ_eq σ_ppi k : ℝ)
    (h1 : is_spike eq_vol μ_eq σ_eq k)
    (h2 : is_spike ppi_vol μ_ppi σ_ppi k) :
    (eq_vol + ppi_vol) / 2 > (μ_eq + μ_ppi) / 2 + k * (σ_eq + σ_ppi) / 2 := by
  unfold is_spike at *
  linarith
\end{lstlisting}
\end{remark}

\subsection{Rolling Mean Bounded by Maximum}

\begin{theorem}[Rolling Mean Bounded by Window Maximum]
\label{thm:mean_bound}
For a finite sequence $v : \{0,\ldots,n\} \to \mathbb{R}$ and an upper
bound $M$ such that $v_i \leq M$ for all $i$:
\[
    \frac{1}{n+1}\sum_{i=0}^{n} v_i \;\leq\; M
\]
\end{theorem}

\begin{proof}
\[
    \sum_{i=0}^{n} v_i \;\leq\; \sum_{i=0}^{n} M \;=\; (n+1)\cdot M
\]
Dividing both sides by $n+1 > 0$ gives the result. \qed
\end{proof}

\begin{remark}
This theorem formalises why a larger rolling window smooths volatility
estimates — the window mean can never exceed the largest value in the
window. For the 90-day window used in APVI, this ensures that a single
extreme day cannot drive the rolling estimate above its own value.
\end{remark}

% ---------------------------------------------------------------
\section{Empirical Interpretation}
% ---------------------------------------------------------------

The following table summarises the correspondence between the formal
theorems and observed market events:

\begin{center}
\begin{tabular}{lll}
\toprule
\textbf{Event} & \textbf{Observed Effect} & \textbf{Relevant Theorem} \\
\midrule
COVID-19 (Mar 2020)        & XAR spike + PPI spike    & Theorem~\ref{thm:apvi_spike} \\
Suez Canal blockage (2021) & XAR spike                & Theorem~\ref{thm:spike_monotone} \\
Ukraine conflict (2022)    & PPI spike                & Theorem~\ref{thm:apvi_bound} \\
General baseline           & APVI $\geq 0$            & Theorem~\ref{thm:apvi_nonneg} \\
\bottomrule
\end{tabular}
\end{center}

% ---------------------------------------------------------------
\section{Conclusion}
% ---------------------------------------------------------------

These six formally verified theorems establish the mathematical
foundations of the APVI model:

\begin{enumerate}
    \item Spike detection is monotone — stricter thresholds are harder to trigger.
    \item Realised volatility is always non-negative.
    \item The composite APVI is always non-negative.
    \item The APVI is bounded above by its most extreme component.
    \item Bilateral spikes in both equity and PPI necessarily produce a composite spike.
    \item Rolling window means are bounded by the window maximum.
\end{enumerate}

These properties together guarantee that the APVI is a well-defined,
interpretable, and mathematically sound risk metric. The full Lean 4
source code is available at:
\begin{center}
\url{https://github.com/Dean-Foulds/lean_projects}
\end{center}

\end{document}
